\scanpage{248}
\margin{B) \uline{Case of punctured surfaces}}
\deleted{\unsure{In short.}}

That said, let us return to the situation of (53), of a topological space $X$, and of an $x\in X$ such that $X$ is a surface at $x$. Then the category \deleted{of \unsure{coverings}}
\[
\underline{\mathrm{Rev}}(D_x^*)
\]
is a Galois category, \inserted{abelian,} whose fundamental group is identified \unsure{canonically} with the orientation module $\overline T_x$ of $X$ at $x$
\[
(56)\qquad \pi_1(D_x^*)\overset{\mathrm{def}}=\pi_1(\underline{\mathrm{Rev}}(D_x^*))\simeq\overline T_x.
\]
This comes from the fact that the germ $D_x^*$ \deleted{\unsure{is} \illegible{1}} \inserted{admits a ``basis''} formed of spaces $D_{x,i}^*$ \deleted{\unsure{isomorphic to}} \inserted{\unsure{arising}} from the \uline{disks} with center $x$, and the $\pi_1$ of the said disks \unsure{form} a constant projective system\ldots

Thus the universal coverings $D_x^{*\sim}$ of $D_x^*$ form a connected groupoid, with fundamental \unsure{groupoid|group} $\deleted{\unsure{\mathbb{Z}}}\,\overline T_x$.

Let us therefore choose such a $D_x^{*\sim}$, and consider the composite morphisms

\scanpage{249}
\mspage{772}
\diagram{\begin{tikzcd}
D_x^{*\sim}\arrow[r] & D_x^*\arrow[r,hook] & X-\{x\}
\end{tikzcd}}
- more generally, let $S$ be a subset of $X$, containing $x$ as an isolated point (for example a discrete subset of $X$), so that we have
\diagram{\begin{tikzcd}
(57)\qquad D_x^*\arrow[r,hook] & U\overset{\mathrm{def}}=X\setminus S
\end{tikzcd}}
and consider the composites
\diagram{\begin{tikzcd}
(58)\qquad D_x^{*\sim}\arrow[r] & D_x^*\arrow[r,hook] & U.
\end{tikzcd}}
If $\underline{\mathrm{Rev}}(U)$ is \unsure{multi-Galois} (for example $U$ \inserted{locally} $\unsure{1|0}$-connected), then we can therefore consider the corresponding fundamental group of $U$
\[
(59)\qquad \pi_1(U,D_x^{*\sim}),
\]
I am going to define a \uline{canonical} \uline{homomorphism}
\diagram{\begin{tikzcd}[column sep=large]
(60)\qquad \unsure{K_x|\kappa_x}:\overline T_x\arrow[r] &
\underbrace{\pi_1(U,D_x^{*\sim})}_{\inserted{\begin{gathered}
\simeq\operatorname{Aut}\text{ of the fiber}\\
\text{functor }\underline{\mathrm{Rev}}(U)\longrightarrow\underline{\mathrm{Rev}}(D_x^{*\sim})\\
\text{on }\underline{\mathrm{Rev}}(U)
\end{gathered}}},
\end{tikzcd}}
we can describe it for example as
\diagram{\begin{tikzcd}
(61)\qquad \pi_1(D_x^*,D_x^{*\sim})\arrow[r]\arrow[d,phantom,"\simeq" description,sloped] & \pi_1(U,D_x^{*\sim})\\
\overline T_x &
\end{tikzcd}}

\scanpage{250}
induced by the inclusion
\diagram{\begin{tikzcd}
D_x^*\arrow[r,hook] & U
\end{tikzcd}}
or again, \deleted{\unsure{by}} in terms of \unsure{automorphisms} of a fiber functor, as the \deleted{\unsure{homomorphism}}operation of $\overline T_x$ on
\[
\Gamma\bigl(D_x^{*\sim},\underbrace{U'\mathbin{\times}_{\deleted{\unsure{D}}U}D_x^{*\sim}}_{\inserted{U'_{D_x^{*\sim}}}}\bigr)
\]
coming from the operations of $\overline T_x$ on $D_x^{*\sim}$.
\diagram{\begin{tikzcd}[column sep=small,row sep=small]
& U'_{D_x^{*\sim}}\arrow[dl]\arrow[dr] &\\
D_x^{*\sim}\arrow[dr] && U'\arrow[dl]\\
& U &
\end{tikzcd}}
\deleted{\illegible{1} \uline{N.B.} \unsure{Each time}}

Suppose that we have already chosen \inserted{separately} a fundamental group \inserted{$\pi_1(U,\widetilde U)=\pi$} of $U$, by the choice of a base point, or, more \unsure{generally}, of a universal covering $\widetilde U$ of $U$.
To compare it with $\pi_1(U,\deleted{\unsure{S}}D_x^{*\sim})$, we shall therefore choose a ``path of \unsure{domination}'' from $D_x^{*\sim}/U$ \unsure{to|:} $\widetilde U/U$, i.e. an isomorphism
\diagram{\begin{tikzcd}
(62)\qquad \widetilde U\arrow[r,"\sim"] & \mathrm{Revuni}(U,D_x^{*\sim})
\end{tikzcd}}
(universal covering of $\deleted{\unsure{\widetilde U}}U$ defined by the fiber functor
\diagram{\begin{tikzcd}
& \underline{\mathrm{Rev}}(U)\arrow[r] & \underline{\mathrm{Rev}}(D_x^{*\sim})\\
\inserted{\unsure{U}}\arrow[ur,bend right] &&
\end{tikzcd}}
\margin{\uline{N.B.} if \deleted{\illegible{1}} $D_x^{*\sim}$ defined by a universal covering $D_{x,i}^{*\sim}$ of a \unsure{punctured} disk $D_{x,i}^*$, and if $\widetilde y\in D_{x,i}^{*\sim}$, $\widetilde z\in\widetilde U$, $y,z$ the images in $U$, then
\[
\deleted{\unsure{D}}\mathrm{Revuni}(U,D_x^{*\sim})\simeq\mathrm{Revuni}(U,y)
\]
and $\widetilde U=\mathrm{Revuni}(U,z)$ : the choice of (62) \unsure{is equivalent to} that of a \unsure{path} from $y$ to $z$ ($\unsure{U}$ \unsure{locally} \unsure{1-connected}).}

\scanpage{251}
\mspage{773}
(the indeterminacy is \unsure{modulo} an \deleted{\unsure{automorphism}} \inserted{element} \deleted{\unsure{inner}} of $\pi$, or of $\pi_1(U,D_x^{*\sim})$) which induces an isomorphism
\diagram{\begin{tikzcd}
(63)\qquad \unsure{\pi_z|\pi_s|\pi}=\pi_1(U,\widetilde U)\arrow[r,leftrightarrow,"\sim"] & \pi_1(U,D_x^{*\sim})
\end{tikzcd}}
(which is changed by an inner automorphism, when we \unsure{change|compose} (62) by an element of $\pi$).

Composing (63) and (60), we therefore find a homomorphism
\diagram{\begin{tikzcd}
(64)\qquad \overline T_x\arrow[r] & \pi=\pi_1(U,\widetilde U),
\end{tikzcd}}
which depends on the choice of (62) -- and which, independently of this choice, defines a well-determined \uline{outer} \uline{homomorphism} -- which is none other than the outer homomorphism
\diagram{\begin{tikzcd}
(65)\qquad \pi_1(D_x^*)\arrow[r]\arrow[d,phantom,"\simeq" description,sloped] & \pi_1(U)\arrow[d,phantom,"\simeq" description,sloped].\\
\overline T_x & \pi
\end{tikzcd}}

\margin{C) \uline{Relation with} \uline{germs of arcs} (\unsure{playing} the role \deleted{\unsure{origin}} of ``base \unsure{points|point}'')}
I am going to describe a convenient way of fixing a universal covering $D_x^{*\sim}$ of $D_x^*$, by a \uline{germ} \uline{of an arc of a curve} \inserted{\uline{of} $X$,} \uline{issuing}

\scanpage{252}
from $x$. In general, let $I=[0,1]$, and $[I]$ be the \unsure{residual} germ of $I$ at $0$. \deleted{If $X$ is a space} \inserted{This is obviously a \uline{1-connected} \uline{germ}.}

\margin{Rather than \uline{\unsure{parametrize}} by germs of \unsure{neighborhoods} of $0$ in $[0,1]$, it would be more convenient for what follows to define a germ of an arc of $X$ issuing from $x$ as a germ of a \unsure{subset} of $X$ at $x$, satisfying \unsure{certain} properties\ldots.}

If $X$ is a topological space and $x\in X$, we call a germ of an arc of a curve \deleted{of} $X$ issuing from $x$, the morphisms
\diagram{\begin{tikzcd}
(66)\qquad [I]\arrow[r,"\gamma"] & \unsure{\mathrm{germ\,rés}}_x(X)\overset{\mathrm{def}}=D_x^*
\end{tikzcd}}
from the germ $[I]$ into the \unsure{residual} germ \inserted{$D_{X,x}^*$} of $X$ at $x$. Such a germ is therefore described by a continuous map from an interval $[0,\varepsilon]$ ($0<\varepsilon\leq1$)
\diagram{\begin{tikzcd}
(67)\qquad [0,\varepsilon]\arrow[r,"{\varphi\,\deleted{\unsure{\gamma}}}"] & X
\end{tikzcd}}
such that we have
\[
(68)\qquad\left\{\begin{array}{l}
\varphi(0)=x\\
\varphi(t)\neq x\quad\text{if }t\in]0,\varepsilon]
\end{array}\right.\quad.
\]
In the case we are interested in ($X$ a surface in a neighborhood of $x$, hence $\underline{\mathrm{Rev}}(D_x^*)$ \deleted{\unsure{Galois}} \inserted{Galois} abelian with group $\overline T_x$), we can use this morphism to define a fiber functor on $D_x^*$
\diagram{\begin{tikzcd}
(69)\qquad \underline{\mathrm{Rev}}(D_x^*)\arrow[r] & \underline{\mathrm{Rev}}([I]) & (\mathrm{Ens})\arrow[l,"\sim"']
\end{tikzcd}}

\scanpage{253}
\mspage{774}

hence a universal covering
\[
\tag{70}
D_x^{*\sim}\overset{\mathrm{d\acute ef}}{=}
\mathrm{Rev\,uni}\bigl(D_x^*,([I],\varphi)\bigr).
\]

Of course, we shall have a factorization
\diagram{\begin{tikzcd}[column sep=large]
[I] \arrow[r] \arrow[rr,bend right=18,"\varphi"'] &
D_x^{*\sim} \arrow[r] & D_x^*\;( \arrow[r,hook] & U)
\end{tikzcd}}

so we have a \unsure{canonical isomorphism}
\[
\tag{71}
\mathrm{Rev\,uni}(U,D_x^{*\sim})\simeq
\mathrm{Rev\,uni}(U,[I]),
\]
where $I$ is regarded as a $1$-connected germ over $U$ by the composite morphism
\diagram{\begin{tikzcd}
{[I]} \arrow[r] & D_x^* \arrow[r] & U.
\end{tikzcd}}

Thus we have
\diagram{\begin{tikzcd}[column sep=large]
{\text{(72)}\quad \pi_1(U,D_x^{*\sim})=\pi_1(U,[I])} &
\overline T_x. \arrow[l,"K_x"']
\end{tikzcd}}

\margin{D) \uline{Relation with differential structure}}
When $X$ is equipped with a differentiable structure in a neighborhood of $x$, we say that a germ of a curve in $X$ issuing from $x$ is \uline{differentiable}, if \unsure{the|the} \deleted{map}\inserted{germ} of a continuous map
\diagram{\begin{tikzcd}
I \arrow[r] & X
\end{tikzcd}}
which defines it is differentiable, \deleted{\unsure{Let us take|For this}}\inserted{and moreover}

\scanpage{254}

the tangent map at the origin, \deleted{\illegible{2} \unsure{leads to} a \unsure{morphism}}
\diagram{\begin{tikzcd}
{\text{(73)}\quad\mathbb R} \arrow[r] & \mathrm{Tang}_x(X),
\end{tikzcd}}
is injective. Then the image of $\mathbb R^+$ is a homogeneous half-line in the tangent space to $X$ at $x$, which we call the \uline{half-tangent} of the germ of an arc of a curve under consideration at $x$.

Returning for a moment to the purely topological case, we are going to consider the set $\unsure{\mathrm{Germarc}|\mathrm{Courarc}}(X,x)$ of germs of arcs of $X$ issuing from $x$, as the set of objects of a groupoid (the local fundamental groupoid of $X$ at $x$), taking as morphisms from $\gamma$ to $\gamma'$ \deleted{\unsure{a|the}} the isomorphisms from \deleted{$\mathrm{Rev\,uni}(D_x^*,\gamma)$} $D_{x,\gamma}^{*\sim}$ to $D_{x,\gamma'}^{*\sim}$
\[
\tag{74}
\mathrm{Isom}(\gamma,\gamma')\overset{\mathrm{d\acute ef}}{=}
\mathrm{Isom}_{D_x^*}(D_{x\gamma}^{*\sim},D_{x,\gamma'}^{*\sim}).
\]
We thus find \inserted{by construction} an equivalence of \deleted{categories (of} connected groupoids
\diagram{\begin{tikzcd}[column sep=large]
{\text{(75)}\quad\uline{\unsure{\mathrm{Germarc}|\mathrm{Courarc}}}(X,x)} \arrow[r] &
\uline{\mathrm{Rev\,uni}}(D_x^*).
\end{tikzcd}}